\documentclass[11pt]{article}
\usepackage[margin=0.95in]{geometry}
\usepackage{amsmath,amssymb,amsthm}
\usepackage{fontspec}
\newfontfamily\directorfont[Path=fonts/]{ProofRushDirectorJP.otf}
\usepackage[hidelinks,pdftitle={The exact prime-difference strata of primitive words},pdfauthor={ProofRush Research Team; research directed by Yoshiaki Otsuka}]{hyperref}
\newtheorem{theorem}{Theorem}
\newtheorem{lemma}[theorem]{Lemma}
\newtheorem{corollary}[theorem]{Corollary}
\newcommand{\N}{\mathbb N_0}
\newcommand{\Z}{\mathbb Z}
\newcommand{\Q}{\mathbb Q}
\newcommand{\supp}{\operatorname{supp}}
\newcommand{\rankset}{\mathcal S}
\title{The exact prime-difference strata of primitive words}
\author{ProofRush Research Team\\\small Research directed by {\directorfont 大束良明} (Yoshiaki Otsuka)}
\date{27 November 2026 \qquad Preprint}
\begin{document}
\maketitle
\begin{abstract}
K\'aszonyi conjectured in 2010 that the intermediate classes in his
classification of primitive binary words have stratified semilinear
exponent sets. We prove that every intermediate class fails this condition,
for all multiplicities of the prime factors of the number of marked
letters. The same obstruction corrects the rank-one assertion of his
Theorem~6. We also determine when a proper-rank class contains an
infinite context-free sublanguage. In the only case where it does, each
such sublanguage has an exponent set that is a finite union of points
and arithmetic progressions, each adding equal amounts in two opposite
positions. A Lean proof uses the binary words themselves to refute the
published universal conjecture at 105 marked letters.
\end{abstract}

\section{The conjecture}
A word is primitive if it is nonempty and is not a proper power. For
distinct letters \(a,b\), let \(Q_n\) consist of the primitive words of
the form
\[
 w(e)=a b^{e_0}a b^{e_1}\cdots a b^{e_{n-1}},\qquad e\in\N^n .
\]
Write \(E(Q_n)=\{e:w(e)\in Q_n\}\). The long-standing question about
primitive words led D\"om\"osi and coauthors to study these fixed
regular slices \cite{domosi}. K\'aszonyi and Katsura proved that
\(Q_n\) is context-free when \(n\) has at most three distinct prime
factors \cite{three}. K\'aszonyi subsequently sought grammars for
subclasses of \(Q_n\) and classified its exponents through prime
differences \cite[pp.~314, 325, 328]{kasonyi}. We study his explicitly
stated conjecture about the individual classes.

Let \(P=\{p_1<\cdots<p_k\}\) be the distinct prime divisors of \(n\).
On cyclic indices in \(\Z/n\Z\) put \(T_p e(i)=e(i+n/p)\). For a
subset \(S\subseteq P\) and integers \(t_p\) coprime to their
respective primes put
\[
  D_{S,t}=\prod_{p\in S}(1-T_p^{t_p}).
\]
This is the alternating box sum in Definitions~4--6 of \cite{kasonyi}.
Definition~4 labels a selected denominator \(q_i\) but prints \(p_i\)
in the coprimality condition. Lemma~2 takes inverses modulo \(q_i\),
which identifies the intended prime. In Definition~6 the numeric
difference is compared to an empty-set symbol. Lemma~11 uses the
numeric comparison with zero.
Write \(D_S\) for the unit multipliers. The statements
\(D_{S,t}e\equiv0\) and \(D_Se\equiv0\) are equivalent. In either
direction each cyclic shift is a power of the shift in the other
direction on that axis. Factoring \(1-U^m\) as
\((1-U)(1+\cdots+U^{m-1})\) proves the assertion with positive
residue representatives. Operators on different axes commute.

Let \(V_\alpha\subseteq\Q^n\) consist of the arrays with every
difference of order \(\alpha+1\) identically zero, and put
\(V_k=\Q^n\). All higher differences of a member of \(V_\alpha\)
also vanish. The exact primitive class \(\rankset_\alpha=E_\alpha\cap
E(Q_n)\) in Definition~10 of \cite{kasonyi} lies in \(V_\alpha\) and
asks for nonvanishing differences at order \(\alpha\). Equation~(57)
can be read as requiring nonvanishing for every prime subset of that
size. Lemma~11(ii) specifies existence at each order up to
\(\alpha\), as does the pattern in Example~2 on p.~329. Our witnesses
have nonvanishing differences for every
nonempty subset of every order through \(\alpha\).

A set of vectors is stratified semilinear if it is a finite union of
linear sets \(b+\sum_j\N v_j\), with \(b,v_j\in\N^n\), in each of
which the periods \(v_j\) have support of size at most two and no two
of their supports cross. Crossing means alternating endpoints
\(i<j<\ell<m\) of the two pairs. The classical bounded-language
characterization applies with these conditions \cite{ginsburg}.

\begin{theorem}\label{main}
For every \(n\) with \(k\ge2\) and \(1\le\alpha<k\), the set
\(\rankset_\alpha\) is not stratified semilinear. In fact it has no
finite union presentation by linear sets with periods supported on at
most two positions, even when the periods may cross. The top class
\(\rankset_k\) is stratified semilinear.
\end{theorem}

The range \(2\le\alpha\le k-1\) answers the conjecture on printed
p.~328. The case \(\alpha=1\) corrects the first-rank clause of
Theorem~6 on p.~325. The formal counterexample below occurs inside
\(Q_{105}\), a language already known to be context-free. The larger
questions about all the \(Q_n\) and about the full primitive-word
language remain open.

\section{Possible periods inside a mixed-difference kernel}
Write \(d=\prod_{p\in P}p\) and \(\theta=n/d\). In each residue class
modulo \(\theta\), express an index uniquely as \(u+\theta z\) for
\(0\le u<\theta\), \(z\in\Z/d\Z\). The Chinese remainder theorem
identifies \(z\) with a point of \(\prod_{p\in P}\Z/p\Z\).
The shift \(n/p\) changes only its \(p\)-coordinate, by the nonzero
residue \(d/p\bmod p\). All conditions defining \(V_\alpha\)
therefore apply independently on each of these \(\theta\) grids.

\begin{lemma}\label{distance}
Let a finite grid have pairwise distinct side lengths
\(2\le m_1<\cdots<m_k\), and \(0\le\alpha\le k\).
Suppose all \((\alpha+1)\)-axis alternating differences of an array
vanish, for every choice of two levels on each selected axis.
The smallest support of a nonzero such array has size
\(d_\alpha=\prod_{j=1}^{k-\alpha}m_j\).
Equality holds precisely for a nonzero constant times the indicator
of a fiber with the largest \(\alpha\) coordinates fixed.
\end{lemma}
\begin{proof}
When \(\alpha=k\), a singleton suffices and describes every minimizer.
When \(\alpha=0\), every array is constant. For \(0<\alpha<k\), slice
along the axis of length \(m_1\).
If two slices differ, their difference has vanishing \(\alpha\)-axis
differences on the remaining grid. Its support is bounded above by
the sum of the supports of the two slices. Induction applied to that
difference bounds the support of the entire array below by
\(\prod_{j=2}^{k-\alpha+1}m_j>d_\alpha\).
If all slices coincide, their common array satisfies the original
order bound on the remaining grid. Induction gives a lower bound
\(m_1\prod_{j=2}^{k-\alpha}m_j=d_\alpha\), and describes every
equality case. The listed indicators depend only on their fixed
coordinates, so they attain the bound.
\end{proof}

The distance value also follows from the minimum distance of a downset
evaluation code \cite[Lemma~2]{downset}. On distinct rational labels
one can interpolate each axis up to degree \(m_i-1\). The monomials of
interest use at most \(\alpha\) variables: expansion about a fixed grid
point by inclusion--exclusion decomposes any array into its anchored
mixed differences, and all summands of higher order vanish. This
identifies the relevant downset. The equality case is also a
specialization of Geil's footprint results \cite[Theorem~7 and
Corollary~9]{geil}: distinct side lengths give a unique allowed
minimum-distance monomial. We have included the induction to make the
application to the prime-difference spaces self-contained.

\begin{corollary}\label{periods}
If \(\alpha<k\), a nonzero vector of \(V_\alpha\) with support of size
at most two exists precisely when \(n\) is even and \(\alpha=k-1\).
In that case it is a scalar multiple of
\(h_i=\mathbf e_i+\mathbf e_{i+n/2}\) for some \(0\le i<n/2\).
\end{corollary}
\begin{proof}
Apply Lemma~\ref{distance} on each nonzero residue class. A
two-position support is possible only when the product of the
\(k-\alpha\) smallest distinct primes is 2. Both positions then lie in
one class, have equal entries, and fill its order-two axis.
\end{proof}

A second proof of Corollary~\ref{periods} works directly in cyclic
indices. Let \(\zeta\) be a primitive \(n\)th root of unity and
\(\widehat v(m)=\sum_i v(i)\zeta^{mi}\). Vanishing of the full
difference forces \(\widehat v(1)=0\), since its multiplier
\(\prod_{p\mid n}(1-\zeta^{-n/p})\) is nonzero. Two distinct powers
of \(\zeta\) can be linearly dependent over \(\mathbb R\) only at
antipodes. Cancellation then requires equal coefficients. If
\(\alpha\le k-2\) and \(n\) is even, choose \(\alpha+1\) odd prime
axes. Their multiplier at \(m=2\) is nonzero, forcing
\(\widehat v(2)=0\), whereas an antipodal pair of weight \(c\ne0\)
has coefficient \(2c\zeta^{2i}\). For \(\alpha=k-1\) an antipodal
pair is killed by \(1-T_2\).

\section{Primitive witnesses and the exact classes}
Primitivity has a convenient cyclic test. A decomposition
\(w(e)=v^m\), \(m\ge2\), starts each copy of \(v\) with \(a\).
If \(v\) has \(j\) occurrences of \(a\), then \(n=mj\) and the gap
list repeats every \(j\) positions. For a prime \(p\mid m\), this gives
\(T_pe=e\). Conversely, \(T_pe=e\) writes \(w(e)\) as \(p\) copies
of its first \(n/p\) blocks. Zero gaps cause no exception.

On the residue class \(u=0\), put \(I_p(z)=1\) if \(z\bmod p=0\)
and \(I_p(z)=0\) otherwise. Define
\[
 F_\alpha(z)=\sum_{\substack{J\subseteq P\\|J|=\alpha}}
                 \prod_{p\in J} I_p(z)
\]
there, and put it equal to zero on the other classes. Each summand
omits an axis of every higher-order difference. At \(z=0\) all
admissible \(S\)-differences with \(1\le |S|\le\alpha\) satisfy
\[
 (D_{S,t}F_\alpha)(0)=\binom{k-|S|}{\alpha-|S|}>0 .
\]
Only terms indexed by \(J\supseteq S\) contribute. In particular every
prime shift fails to preserve \(F_\alpha\). The distinct vectors
\(mF_\alpha\) for \(m\ge1\) belong to the required primitive class.

\begin{proof}[Proof of Theorem~\ref{main}]
Every component in an exact finite-union representation is contained
in \(V_\alpha\). Both its base \(b\) and \(b+v\) lie in this vector
space for each period \(v\), so \(v\in V_\alpha\).
Outside the exceptional case of Corollary~\ref{periods}, all periods
are zero and a finite union cannot contain the infinite family above.
In the exceptional case \(D_{\{2\}}v=0\) for every allowed period.
The number \((D_{\{2\}}e)(0)\) is fixed on each component and takes
only finitely many values in their union. On \(mF_\alpha\) it takes
the distinct values \(m\binom{k-1}{\alpha-1}\).

At order \(k\), a nonzero full difference implies that each prime
period is violated. Thus the top set is the union over the finitely
many base positions of a nonzero signed linear comparison of the gap
lengths. Each such comparison has a one-counter test: while scanning
the fixed blocks, add the signed weight of each letter to a counter
and accept for a nonzero final balance. A stack of absolute value
with one sign in the finite control implements this test. Finite
union and the bounded-language characterization give the assertion.
\end{proof}

The local comparisons in the proof of Theorem~6 of \cite{kasonyi}
explain the discrepancy. For example, its (79) retains two unequal
gap comparisons guaranteeing primitivity but omits the equalities
imposing rank one. At \(n=15\), take \(e(i)=\mathbf1_{3\mid i}+
\mathbf1_{5\mid i}\). The pairs \((0,5)\) and \((7,10)\) are
noncrossing and have unequal values. Adding one at index \(1\)
preserves both comparisons and makes \(D_{\{3,5\}}e(1)=1\).
Such comparisons can cover a class by larger
primitive sets. Equality with the exact class needs the missing
containment.

\section{Which strata contain an infinite context-free language?}
For fixed \(n\), \(w(E)\) is context-free exactly when \(E\) is
stratified semilinear. Indeed, one may tag each gap by its position
using finite control, obtaining the bounded language with \(n\) distinct
gap letters. The inverse relation inserts the \(n\) markers and erases
the tags, including when a gap has length zero. Both operations are
finite-state transductions. Apply the letter-bounded theorem
\cite[Theorem~2.1]{ginsburg} to the tagged gap lengths. Ginsburg and
Spanier also gave a sparse-period criterion for existence of an
infinite context-free subset of a semilinear bounded language
\cite[Theorem~2.3 and Example~(3) on p.~294]{ginsburg}.
Each present exponent set is semilinear in the ordinary sense:
for fixed \(n\) its defining equations, disequations and the
prime-period tests form a finite Boolean combination of linear
conditions. We identify which sparse periods actually occur for this
family and determine the permitted sublanguages.

\begin{theorem}\label{immune}
Let \(1\le\alpha<k\). If \(n\) is odd or \(\alpha\le k-2\), every
context-free sublanguage of \(w(\rankset_\alpha)\) is finite.
If \(n\) is even and \(\alpha=k-1\), the class has an infinite
context-free sublanguage. A sublanguage of the class is context-free
if and only if its exponent set is a finite union of points and
arithmetic progressions whose nonzero steps have the form \(c h_i\)
for \(c\in\mathbb N\).
\end{theorem}
\begin{proof}
Apply Corollary~\ref{periods} to any stratified presentation of the
sublanguage. In the even exceptional case two distinct antipodal
pairs cross: for \(0\le i<j<n/2\), their order is
\(i<j<i+n/2<j+n/2\). Thus within a single component all nonzero
periods lie on the same ray. Its scalar finitely generated submonoid
of \(\N\) is a finite union of arithmetic progressions or a singleton.
Indeed, given a positive generator \(c\), take the least represented
element in each occupied residue class modulo \(c\) and add
\(0,c,2c,\ldots\). Conversely each stated progression is already
stratified linear.

To see existence, take \(F=F_{k-1}\) and fix one \(h_i\).
Every \(F+t h_i\) has zero full difference. For each nonempty \(S\)
of size at most \(k-1\), its unit-scale difference at zero is affine
in \(t\), with a strictly positive constant term. It can vanish for
at most one \(t\). The finitely many exceptions over all \(S\) are
avoided by all sufficiently large \(t\). The remaining arithmetic
progression lies in the exact primitive class.
\end{proof}

In the standard terminology \cite{immune}, this proves CFL-immunity
of the infinite exact classes outside the even exceptional case.
It rules out generating any infinite portion of those classes by
a context-free grammar constrained to stay inside the given class.
Constructions for the whole \(Q_n\) may use grammars that cross ranks.

\section{A binary-word certificate}
For \(n=105\), \(\alpha=2\), let \(I_p(i)=\mathbf1_{p\mid i}\). Use
\[
  f=I_3I_5+2I_3I_7+2I_5I_7,\qquad e_t=(105t+1)f\quad(t\in\N).
\]
The pair differences at zero have coefficients \(1,2,2\) for all
admissible multipliers. All triple differences vanish. The sum of
the gaps in \(f\) is \(7+2\cdot5+2\cdot3=23\). Hence \(w(e_t)\)
has \(105\) letters \(a\) and \(23(105t+1)\) letters \(b\); the two
counts are coprime. Every proper repetition exponent would divide
both. The zeroth gaps \(5(105t+1)\) grow without bound.

The supplied Lean modules define binary words as lists, encode each
gap literally, quantify over arbitrary integer scale multipliers, and
define every period and the noncrossing condition. They exclude finite
stratified presentations of the exact set at these parameters, both for
universal and existential readings of order-two nonvanishing. The
all-\(n\) grammar and immunity conclusions use the arguments above
and the classical bounded-language theorem. The research archive
identifies the precise formal targets and supplies a clean replay.

\paragraph{Research and review.}
Yoshiaki Otsuka set the research objective and acceptance criteria.
OpenAI Codex agents developed the argument and carried out the
formalization, literature search and independent mathematical checks.
External specialist review is pending. The source question has no
identified prize of its own.

\begin{thebibliography}{9}
\bibitem{domosi}
P. D\"om\"osi, S. Horv\'ath, M. Ito, L. K\'aszonyi and M. Katsura,
Formal languages consisting of primitive words, \emph{Fundamentals
of Computation Theory}, LNCS 710 (1993), 194--203.
\href{https://doi.org/10.1007/3-540-57163-9_15}{doi:10.1007/3-540-57163-9\_15}.
\bibitem{three}
L. K\'aszonyi and M. Katsura, On the context-freeness of a class of
primitive words, \emph{Publ. Math. Debrecen} 51 (1997), 1--11.
\href{https://doi.org/10.5486/PMD.1997.1707}{doi:10.5486/PMD.1997.1707}.
\bibitem{kasonyi}
L. K\'aszonyi, Semiperiodic vectors and the context-freeness of
\(Q_n=Q\cap(ab^*)^n\), \emph{Publ. Math. Debrecen} 77 (2010), 313--330.
\href{https://doi.org/10.5486/PMD.2010.4680}{doi:10.5486/PMD.2010.4680}.
\bibitem{ginsburg}
S. Ginsburg and E. H. Spanier, Semigroups, Presburger formulas, and
languages, \emph{Pacific J. Math.} 16 (1966), 285--296.
\href{https://doi.org/10.2140/pjm.1966.16.285}{doi:10.2140/pjm.1966.16.285}.
\bibitem{downset}
S. Srinivasan, U. Tripathi and S. Venkitesh, Decoding downset codes
over a finite grid (2019). \href{https://arxiv.org/abs/1908.07215}{arXiv:1908.07215}.
\bibitem{geil}
O. Geil, On multivariate polynomials with many roots over a finite
grid. \href{https://arxiv.org/abs/1909.06613}{arXiv:1909.06613}.
\bibitem{immune}
T. Yamakami, Immunity and pseudorandomness of context-free languages.
\href{https://arxiv.org/abs/0902.0261}{arXiv:0902.0261}.
\end{thebibliography}
\end{document}
